\clearpage
\section{Convergence to the stationary states}\label{sec:convergence}

\noindent\motiv{\reve{A reversible dynamics does not bring the state of the chain to a limit in norm
unless the state is stationary, since the distance between the states at consecutive times does not
change with time. What it can do is spread the correlations of any finite region over ever larger
distances, so that every expectation value settles to a limit. For the example system in the regime
this happens for every Gaussian initial state with zero displacement whose covariance is given by
cell data. For reflection-symmetric data the symbol of the limit is, wave number by wave number, of
the form found in Theorem~\ref{thm:es-stationary}.
After $t$ steps the symbol $\hat L(k)^t$ is conjugate to the rotation by the angle $t\theta(k)$, so that
the evolved symbol of the covariance consists of a part independent of $t$ and parts proportional to
$\cos(2t\theta(k))$ and $\sin(2t\theta(k))$. The cells are the Fourier coefficients of this symbol.
Because $\theta(k)$ varies with $k$, the Fourier coefficients of the oscillating parts tend to zero,
and the limit is the average of the symbol over the rotation, which is invariant under
$S\mapsto\hat L(k)^{\mathsf T}S\hat L(k)$.}}

\medskip
\noindent\motiv{\reve{As in the previous section, $(\varphi,\lambda)$ lies in the regime of Definition~\ref{def:es-regime}, and $\theta(k)$ is
the rotation angle of Definition~\ref{def:es-alpha}. The real matrix $N(k)$ is defined by
$\hat L(k)=\cos\theta(k)\,1-\sin\theta(k)\,N(k)$ and satisfies $N(k)^2=-1$, and $\hat\varepsilon(k)$ is the
invariant form of Definition~\ref{def:es-eps}.}}

\subsection{Oscillatory integrals}

\noindent\motiv{\reve{The Riemann--Lebesgue lemma states that $\int f(u)e^{\ii ru}\,du$ tends to zero as
$r\to\infty$ for every integrable $f$. The integrals that arise here carry the phase $r\theta(k)$, with
$\theta$ the rotation angle, instead of $ru$. Since $\theta$ is strictly monotone on each half of the
circle, a change of variables brings them to the standard form.}}

\begin{lemma}[Oscillatory integrals along a monotone phase]\label{lem:cv-rl}
\reve{Let $a\le b$, let $g$ be strictly monotone on $[a,b]$ and differentiable on $(a,b)$ with
$g'\neq0$ there, and let $f$ be integrable on $[a,b]$. Then
\[\int_a^bf(k)\,e^{\ii rg(k)}\,dk\longrightarrow0\qquad(r\to+\infty).\]}
\end{lemma}

\begin{proof}
\reve{The endpoints have measure zero, so that we may take the integral over the open interval
$(a,b)$, on which $g$ is injective with a nonvanishing derivative. Its image $I=g((a,b))$
is measurable, and $g$ has an inverse $g^{-1}\colon I\to(a,b)$. The change-of-variables formula
for an injective differentiable map gives
\[\int_IG(u)\,du=\int_a^b|g'(k)|\,G(g(k))\,dk\]
for every measurable function $G$, with $G$ integrable on $I$ exactly when $|g'|\,(G\circ g)$ is
integrable on $(a,b)$. For $G(u)=H(u)e^{\ii ru}$ with $H=(f/|g'|)\circ g^{-1}$ the right side is the
integral in question, and $|g'|\,(H\circ g)=f$ is integrable, so that $H$ is integrable
on $I$. Hence
\[\int_a^bf(k)\,e^{\ii rg(k)}\,dk=\int_\R\mathbf 1_I(u)H(u)\,e^{\ii ru}\,du,\]
the integral of the integrable function $\mathbf 1_IH$ against $e^{\ii ru}$, which tends to zero by
the Riemann--Lebesgue lemma.}
\end{proof}

\noindent\reve{When $g'$ tends to zero at an endpoint, as it will for the rotation angle $\theta$ at
$k=0$ and $k=\pm\pi$, the function $H$ may be unbounded there. The proof requires only its
integrability, which the change of variables supplies, so that no separate estimate near the critical
points of $g$ is needed.}

\begin{corollary}[Oscillatory integrals of the example]\label{cor:cv-rl}
\reve{Let $(\varphi,\lambda)$ lie in the regime, and let $f\colon\R\to\C$ be continuous. Then,
for either choice of sign,
\[\int_{-\pi}^{\pi}f(k)\,e^{\pm2\ii t\theta(k)}\,dk\longrightarrow0\qquad(t\to\infty,\ t\in\mathbb N).\]}
\end{corollary}

\begin{proof}
\reve{Split the integral at $k=0$. By Proposition~\ref{prop:es-alpha}, $\theta$ is differentiable with
$\theta'(k)=\lambda\sin\varphi\sin k/\sin\theta(k)$, which vanishes nowhere on $(0,\pi)$ or on
$(-\pi,0)$, because $\lambda\sin\varphi\neq0$ in the regime; on $[0,\pi]$ it is strictly monotone, and on
$[-\pi,0]$ it is strictly monotone in the opposite sense, being even. Lemma~\ref{lem:cv-rl} with $r=2t$,
applied on each half to $\theta$ for the upper sign and to $-\theta$ for the lower, gives the claim.}
\end{proof}

\subsection{The limit of the covariance}

\noindent\motiv{\reve{The evolved cell data $(L_*)^t\gamma$ have the symbol
$(\hat L(k)^t)^{\mathsf T}\hat\gamma(k)\hat L(k)^t$, by Proposition~\ref{prop:es-iterate} and the
reflection symmetry of $\hat L$ (Proposition~\ref{prop:es-symbol}). The power $\hat L^t$ is the
rotation $\cos(t\theta)\,1-\sin(t\theta)\,N$, and the quadratic expression in it separates into a part
that does not depend on $t$ and parts at twice the angle.}}

\begin{lemma}[Splitting of the evolved symbol]\label{lem:cv-split}
\reve{Let $(\varphi,\lambda)$ lie in the regime, and let $A$ be a complex $2\times2$ matrix. At each wave
number $k$, with $\theta=\theta(k)$ and $N=N(k)$, and for every $t\in\mathbb N$,
\[(\hat L(k)^t)^{\mathsf T}A\,\hat L(k)^t=C+\cos(2t\theta)\,A_c+\sin(2t\theta)\,A_s,\]
where $C=\frac12(A+N^{\mathsf T}AN)$, $A_c=\frac12(A-N^{\mathsf T}AN)$ and
$A_s=-\frac12(N^{\mathsf T}A+AN)$.}
\end{lemma}

\begin{proof}
\reve{By Proposition~\ref{prop:es-projectors}(c) and $P_k-\bar P_k=\ii N$, the power is
$\hat L^t=\cos(t\theta)\,1-\sin(t\theta)\,N$, so that
\[(\hat L^t)^{\mathsf T}A\hat L^t=\cos^2(t\theta)\,A-\cos(t\theta)\sin(t\theta)\,(N^{\mathsf
T}A+AN)+\sin^2(t\theta)\,N^{\mathsf T}AN,\]
and the double-angle formulas $\cos^2x=\frac12(1+\cos2x)$, $\sin^2x=\frac12(1-\cos2x)$ and
$\cos x\sin x=\frac12\sin2x$ give the stated form.}
\end{proof}

\noindent\motiv{\reve{The part $C$ is the average of $(\hat L^t)^{\mathsf T}A\hat L^t$ over a full
rotation.}}

\begin{definition}[Limit symbol and limit cells]\label{def:cv-limit}
\reve{Let $(\varphi,\lambda)$ lie in the regime. For cell data $(\gamma_x)$ with symbol $\hat\gamma$,
the \emph{limit symbol} is
\[C(k)=\tfrac12\bigl(\hat\gamma(k)+N(k)^{\mathsf T}\hat\gamma(k)N(k)\bigr),\]
and the \emph{limit cells} are its Fourier coefficients,
\[\gamma_{\infty,x}=\frac1{2\pi}\int_{-\pi}^{\pi}e^{\ii kx}\,C(k)\,dk\qquad(x\in\Z).\]}
\end{definition}

\begin{lemma}[The limit cells]\label{lem:cv-limit}
\reve{Let $(\varphi,\lambda)$ lie in the regime. For cell data $(\gamma_x)$, every $\gamma_{\infty,x}$
is a real matrix, and
$\gamma_{\infty,-x}=\gamma_{\infty,x}^{\mathsf T}$.}
\end{lemma}

\begin{proof}
\reve{The matrix $N(k)$ is real and even in $k$, and $\hat\gamma(-k)=\overline{\hat\gamma(k)}$ by
Lemma~\ref{lem:gs-cells}, so that $C(-k)=\overline{C(k)}$. Conjugating the defining integral and
substituting $k\mapsto-k$ gives
$\overline{\gamma_{\infty,x}}=\frac1{2\pi}\int e^{-\ii kx}C(-k)\,dk=\gamma_{\infty,x}$. Since
$\hat\gamma(k)$ is Hermitian, so is $C(k)$, and the same substitution gives
$\gamma_{\infty,x}^{\mathsf T}=\frac1{2\pi}\int e^{\ii kx}\,\overline{C(k)}\,dk
=\frac1{2\pi}\int e^{-\ii kx}C(k)\,dk=\gamma_{\infty,-x}$.}
\end{proof}

\begin{remark}[Summability of the limit cells]\label{rem:cv-summable}
\reve{Let $(\varphi,\lambda)$ lie in the regime and let $(\gamma_x)$ be cell data. Then
$(\gamma_{\infty,x})$ are cell data, and their symbol is $C$.}
\end{remark}

\begin{proof}
\reve{The $2\pi$-periodic functions $F(k)=\sum_xc_xe^{-\ii kx}$ with $\sum_x|c_x|<\infty$ form an algebra
under pointwise multiplication, since the coefficients of a product are the convolution of the
coefficients, and the convolution of two summable sequences is summable. The entries of $\hat\gamma$
belong to it by the definition of cell data. The function $\tau$ is a trigonometric polynomial with
$|\tau|<1$, so that $1-\tau^2$ is a positive trigonometric polynomial; its square root $\sin\theta$ and
the reciprocal $1/\sin\theta$ extend holomorphically to a strip around the real axis, and the entries
of $N=(\tau\,1-\hat L)/\sin\theta$ are therefore $2\pi$-periodic functions holomorphic in a strip, whose
Fourier coefficients decay exponentially. Hence the entries of $C$ lie in the algebra, $C(k)=\sum_x
c_xe^{-\ii kx}$ with summable matrix coefficients $c_x$, and integrating term by term gives
$\gamma_{\infty,x}=c_x$. By Lemma~\ref{lem:cv-limit} these are real with $\gamma_{\infty,-x}=
\gamma_{\infty,x}^{\mathsf T}$.}
\end{proof}

\begin{theorem}[Convergence of the covariance]\label{thm:cv-converge}
\reve{Let $(\varphi,\lambda)$ lie in the regime and let $(\gamma_x)$ be cell data. Then
for every $x\in\Z$,
\[\lim_{t\to\infty}\bigl((L_*)^t\gamma\bigr)_x=\gamma_{\infty,x}.\]}
\end{theorem}

\begin{proof}
\reve{The evolved data $(L_*)^t\gamma$ are again cell data by Proposition~\ref{prop:es-iterate}(b), so
that by Fourier inversion (Proposition~\ref{prop:es-inversion}) and Lemma~\ref{lem:cv-split} with $A=\hat\gamma(k)$,
\[\bigl((L_*)^t\gamma\bigr)_x=\gamma_{\infty,x}+\frac1{2\pi}\int_{-\pi}^{\pi}e^{\ii kx}\bigl(\cos(2t\theta(k))
A_c(k)+\sin(2t\theta(k))A_s(k)\bigr)\,dk .\]
The symbol $\hat\gamma$ is continuous by Lemma~\ref{lem:gs-cells}, and $N$ is continuous in the
regime, since $\theta$ is and $\sin\theta>0$; hence the entries of $A_c$ and $A_s$ are continuous.
Writing the cosine and the sine through $e^{\pm2\ii t\theta(k)}$ expresses each entry of the integral
as a combination of four integrals of the form of Corollary~\ref{cor:cv-rl}, with $f(k)$ equal to
$e^{\ii kx}$ times an entry of $A_c(k)$ or of $A_s(k)$, and each of them tends to zero.}
\end{proof}

\noindent\motiv{\reve{For $\lambda=0$ every wave number is rotated by the same matrix, and the
covariance of a state that is not invariant under this rotation need not converge.}}

\begin{proposition}[No convergence without coupling]\label{prop:cv-nocoupling}
\reve{For $\varphi=\pi/2$, $\lambda=0$ and the cell data $\gamma_0=\operatorname{diag}(2,\frac12)$,
$\gamma_x=0$ for $x\neq0$, the entry $\bigl((L_*)^t\gamma\bigr)_{0,11}$ equals $2$ for even $t$ and
$\frac12$ for odd $t$; in particular it does not converge.}
\end{proposition}

\begin{proof}
\reve{At $\lambda=0$ the symbol is the constant matrix $R=\begin{psmallmatrix}0&-1\\1&0\end{psmallmatrix}$,
and $\hat\gamma\equiv\gamma_0$, so that the evolved symbol is the constant $(R^t)^{\mathsf T}\gamma_0R^t$
and, by Fourier inversion, equals the cell at the origin. Since $R^2=-1$, this matrix is $\gamma_0$ for
even $t$ and $R^{\mathsf T}\gamma_0R=\operatorname{diag}(\frac12,2)$ for odd $t$.}
\end{proof}

\noindent\motiv{\reve{The hypothesis
$\hat\gamma(k)+\ii\Omega\ge0$ below holds at every $k$ whenever the covariance form of the cell data
satisfies the state condition, by Remark~\ref{rem:gs-fourier-general}. In particular, by
Theorem~\ref{thm:gs-necessity}, it holds for the covariance of a Gaussian state.}}

\begin{proposition}[The limit is stationary]\label{prop:cv-stationary}
\reve{Let $(\varphi,\lambda)$ lie in the regime, let $(\gamma_x)$ be cell data and $k\in\R$.
\begin{enumerate}[label=(\alph*)]
\item $\hat L(k)^{\mathsf T}C(k)\hat L(k)=C(k)$.
\item If $\hat\gamma(k)+\ii\Omega\ge0$, then $C(k)+\ii\Omega\ge0$.
\item If $\gamma_{-x}=\gamma_x$ for every $x$ and
$\hat\gamma(k)+\ii\Omega\ge0$, then $C(k)=\eta\,\hat\varepsilon(k)$ for some real $\eta\ge1$.
\end{enumerate}}
\end{proposition}

\begin{proof}\leavevmode
\reve{\begin{enumerate}[label=(\alph*)]
\item Write $A=\hat\gamma(k)$ and $\theta=\theta(k)$, so that $\hat L=\cos\theta\,1-\sin\theta\,N$. Since
$(N^{\mathsf T})^2=(N^2)^{\mathsf T}=-1$, one has $N^{\mathsf T}CN=\frac12(N^{\mathsf T}AN+A)=C$ and
$N^{\mathsf T}C+CN=\frac12\bigl(N^{\mathsf T}A-AN+AN-N^{\mathsf T}A\bigr)=0$, and therefore
$\hat L^{\mathsf T}C\hat L=\cos^2\theta\,C-\cos\theta\sin\theta\,(N^{\mathsf T}C+CN)+\sin^2\theta\,
N^{\mathsf T}CN=C$.
\item $N$ is real with $N^2=-1$, hence $\operatorname{tr}N=0$ and $\det N=1$, and
$N^{\mathsf T}\Omega N=(\det N)\,\Omega=\Omega$ by Lemma~\ref{lem:qca-2x2}. Therefore
$C+\ii\Omega=\frac12\bigl((\hat\gamma+\ii\Omega)+N^{\mathsf T}(\hat\gamma+\ii\Omega)N\bigr)$, and
$N^{\mathsf T}=N^*$ because $N$ is real, so that both terms are positive semidefinite.
\item By Lemma~\ref{lem:es-refl}, $\hat\gamma(k)$ is real symmetric, hence so is $C(k)$. By (a) and
Corollary~\ref{cor:es-coeff}, $C(k)=\eta\,\hat\varepsilon(k)$ with $\eta=h_k(C(k))$, and $\eta\ge1$ by
(b) and Proposition~\ref{prop:es-threshold}.
\end{enumerate}}
\end{proof}

\begin{example}[The vacuum in the running example]\label{ex:cv-running}
\reve{For $\varphi=\pi/2$ and $0<|\lambda|<1$, the vacuum of Example~\ref{ex:gs-vacuum} has
$\hat\gamma\equiv1$ and, by Example~\ref{ex:es-running-stat}, is not stationary. With
$\sin\theta(k)=\sqrt{1-\lambda^2\cos^2k}$ one finds
$N(k)=(\sin\theta(k))^{-1}\begin{psmallmatrix}-\lambda\cos k&1\\-1&\lambda\cos k\end{psmallmatrix}$, hence
\[N^{\mathsf T}N=\frac1{\sin^2\theta}\begin{pmatrix}1+\lambda^2\cos^2k&-2\lambda\cos k\\-2\lambda\cos
k&1+\lambda^2\cos^2k\end{pmatrix},\qquad
C(k)=\frac1{\sin^2\theta(k)}\begin{pmatrix}1&-\lambda\cos k\\-\lambda\cos k&1\end{pmatrix}.\]
This is $\eta(k)\hat\varepsilon(k)$ with $\eta(k)=(1-\lambda^2\cos^2k)^{-1/2}\ge1$, with excess noise wherever $\cos k\neq0$. By
Remark~\ref{rem:cv-summable} the limit cells are cell data with symbol $C$; they are reflection
symmetric, since $C$ is real and even, and $\eta$ is continuous and even with $\eta\ge1$, so that by
Theorem~\ref{thm:es-stationary} the cells of the vacuum converge to those of a stationary covariance.
Here the entries of $C$ are themselves holomorphic in a strip, and the limit cells decay
exponentially.}
\end{example}

\subsection{The limit state}

\noindent\motiv{\reve{The convergence of the covariance cells yields the convergence of the
characteristic functions, since the covariance form on a vector of finite support involves only
finitely many cells. To pass from characteristic functions to arbitrary observables we use that the
Weyl operators span a dense subspace, together with a bound on states that is uniform in the state.}}

\medskip
\noindent\reve{Let $(\gamma_x)$ be cell data with covariance form $\gamma$, and let $\omega$ be a given
Gaussian state with covariance form $\gamma$ and displacement zero. Its evolution by $t$ steps,
$\omega_t=\omega\circ\alpha_{L^t}$, is by Proposition~\ref{prop:gs-evolve} a Gaussian state with
covariance form $\gamma(L^t\cdot,L^t\cdot)$. The limit cells are cell data by
Remark~\ref{rem:cv-summable}, and we denote their covariance form by $\gamma_\infty$.}

\begin{lemma}[Norm of a state]\label{lem:cv-norm}
\reve{For every state $\omega$ on $\CCR(V,\sigma)$ and every $a\in\CCR(V,\sigma)$,
$|\omega(a)|\le\|a\|$.}
\end{lemma}

\begin{proof}
\reve{The form $(a,b)\mapsto\omega(a^*b)$ is positive semidefinite, and the Cauchy--Schwarz inequality
gives $|\omega(a)|^2=|\omega(1^*a)|^2\le\omega(1)\,\omega(a^*a)=\omega(a^*a)$. The element $a^*a$ is
positive in the C*-algebra $\CCR(V,\sigma)$, with spectrum in $[0,\|a^*a\|]=[0,\|a\|^2]$, so that
$\|a\|^2\,1-a^*a$ has non-negative spectrum and is the square of its self-adjoint square root; positivity
of $\omega$ then gives $\omega(a^*a)\le\|a\|^2$.}
\end{proof}

\begin{theorem}[Convergence of the states]\label{thm:cv-weak}
\reve{Let $(\varphi,\lambda)$ lie in the regime, and let $\omega$, $\omega_t$ and $\gamma_\infty$ be as
above.
\begin{enumerate}[label=(\alph*)]
\item For every $u\in V$, $\chi_{\omega_t}(u)\to\exp\bigl(-\frac14\gamma_\infty(u,u)\bigr)$.
\item There is a Gaussian state $\omega_\infty$ with covariance form $\gamma_\infty$ and displacement
zero such that $\omega_t(a)\to\omega_\infty(a)$ for every $a\in\CCR(V,\sigma)$, that is, $\omega_t\to
\omega_\infty$ in the weak-$*$ topology; $\omega_\infty$ is the only state with this characteristic
function.
\end{enumerate}}
\end{theorem}

\begin{proof}\leavevmode
\reve{\begin{enumerate}[label=(\alph*)]
\item Since the displacement vanishes,
$\chi_{\omega_t}(u)=\exp\bigl(-\frac14\gamma(L^tu,L^tu)\bigr)$, and by Theorem~\ref{thm:gs-evolution} and
Proposition~\ref{prop:es-iterate}, $\gamma(L^tu,L^tu)=\sum_{x,y}u(x)^{\mathsf T}\bigl((L_*)^t\gamma\bigr)_{x-y}\,u(y)$. The sum is finite, each cell converges by Theorem~\ref{thm:cv-converge}, and
the exponential is continuous.
\item On a finite combination $b=\sum_jz_jW(u_j)$ the sequence $\omega_t(b)=\sum_jz_j\chi_{\omega_t}(u_j)$
converges by (a). The span of the Weyl operators is dense in $\CCR(V,\sigma)$, which is its closure; for
arbitrary $a$ and $\zeta>0$ choose such a $b$ with $\|a-b\|<\zeta$. By
Lemma~\ref{lem:cv-norm}, $|\omega_t(a)-\omega_{t'}(a)|\le2\zeta+|\omega_t(b)-\omega_{t'}(b)|$, so that
$(\omega_t(a))_t$ is a Cauchy sequence. Its limit $\omega_\infty(a)$ is linear in $a$ and satisfies
$|\omega_\infty(a)|\le\|a\|$, so that it is continuous, and $\omega_\infty(1)=1$ and
$\omega_\infty(a^*a)\ge0$ hold in the limit because the non-negative reals form a closed subset of $\C$.
By (a), $\chi_{\omega_\infty}(u)=\exp\bigl(-\frac14\gamma_\infty(u,u)\bigr)$, and $\gamma_\infty$ is
symmetric by Lemma~\ref{lem:cv-limit}; hence $\omega_\infty$ is Gaussian with covariance form
$\gamma_\infty$ and displacement zero, and it is unique by Proposition~\ref{prop:gs-pullback}(a).
\end{enumerate}}
\end{proof}

\noindent\reve{Since $\omega$ is a Gaussian state, $\hat\gamma(k)+\ii\Omega\ge0$ for every $k$ by
Theorem~\ref{thm:gs-necessity} and Remark~\ref{rem:gs-fourier-general}. By
Proposition~\ref{prop:cv-stationary} the symbol $C$ of the limit therefore satisfies
$\hat L(k)^{\mathsf T}C(k)\hat L(k)=C(k)$ and $C(k)+\ii\Omega\ge0$ at every wave number, and by
Remark~\ref{rem:cv-summable} and Fourier inversion the limit cells are stationary. For
reflection-symmetric data $C(k)=\eta(k)\hat\varepsilon(k)$ with $\eta(k)\ge1$, so that the limit is one of
the stationary covariances of Theorem~\ref{thm:es-stationary}.}

\subsection{Notes and references}\label{sec:cv-notes}

\begin{enumerate}
\item \reve{\emph{Convergence.} Theorem~\ref{thm:cv-converge} is Proposition~4.7
of~\cite[p.~87]{krueger}.
The thesis assumes a pure initial state of finite correlation length. It treats the oscillatory
integrals by a stationary-phase expansion around the extrema of the rotation angle, written as an
equality, without showing that the contributions near the extrema
vanish~\cite[pp.~85--86]{krueger}. Here we assume only summability of the initial cells, and the
change of variables of Lemma~\ref{lem:cv-rl} reduces the decay to the Riemann--Lebesgue lemma. The
thesis does not state the condition $\lambda\neq0$ explicitly; Proposition~\ref{prop:cv-nocoupling} shows that the
statement fails without it. The thesis remarks that the first moments do not converge unless they
vanish~\cite[p.~86]{krueger}. This holds for translation-invariant displacements, which the dynamics
rotates at the single wave number $k=0$. For a displacement supported on finitely many sites one
expects decay to zero by the same dephasing, now with the phases $e^{\pm\ii t\theta(k)}$, but this is
not proved here.}
\item \reve{\emph{The limit.} Proposition~\ref{prop:cv-stationary}(c) identifies the limit symbol with
the form of Theorem~4.8 of~\cite{krueger}, wave number by wave number. The thesis concludes that the
limit is a stationary state in the sense of cell data, on the grounds that $\hat\varepsilon$ is
continuous~\cite[p.~89]{krueger}. Continuity of a symbol does not by itself imply summability of its
Fourier coefficients. The conclusion holds nevertheless, by the real-analyticity of $N$ and the
algebra property used in Remark~\ref{rem:cv-summable}.}
\item \reve{\emph{The limit state.} Theorem~\ref{thm:cv-weak} is the first part of Theorem~4.9
of~\cite{krueger}, the convergence of expectation values. The second part, convergence of the reduced
density matrices of finite regions in trace norm, is not treated in these notes.
Lemma~\ref{lem:cv-norm} is the standard bound $\|\omega\|=\omega(1)$ for positive functionals on a
unital C*-algebra.}
\end{enumerate}
